% ============================================================
% f2014_shandong_science_q11_flowchart.tex
% 2014 山东卷（理科）第 11 题：程序框图
% 来源 PDF：普通高考/2014/2014山东理.pdf
% 原页：第 1 页（A3 横排）中栏，印面 907，第 11 题题干下方
% 结构：开始 → 输入 x → n=0 → 判断 x^2-4x+3≤0；
%       是：x=x+1 → n=n+1 → 返回判断；否：输出 n → 结束。
% ============================================================
\begin{tikzpicture}[>=Stealth, line width=0.7pt, font=\small,
  x=1cm, y=1cm, line cap=round, line join=round]
  % 开始
  \draw[rounded corners=6pt] (-0.95,0.42) rectangle (0.95,-0.42);
  \node at (0,0) {开始};
  \draw[->] (0,-0.42) -- (0,-0.87);
  % 输入 x
  \draw (-1.45,-0.87) -- (1.45,-0.87) -- (1.25,-1.72) -- (-1.65,-1.72) -- cycle;
  \node at (0,-1.30) {输入 $x$};
  \draw[->] (0,-1.72) -- (0,-2.16);
  % n=0
  \draw (-1.10,-2.16) rectangle (1.10,-2.96);
  \node at (0,-2.56) {$n=0$};
  \draw[->] (0,-2.96) -- (0,-3.45);
  % 判断框
  \draw (0,-3.45) -- (2.55,-4.55) -- (0,-5.65) -- (-2.55,-4.55) -- cycle;
  \node at (0,-4.55) {$x^2-4x+3\leqslant 0$};
  % 是：向下进入循环体
  \draw[->] (0,-5.65) -- (0,-6.10);
  \node at (0.45,-5.90) {是};
  \draw (-1.10,-6.10) rectangle (1.10,-6.90);
  \node at (0,-6.50) {$x=x+1$};
  \draw[->] (0,-6.90) -- (0,-7.35);
  \draw (-1.10,-7.35) rectangle (1.10,-8.15);
  \node at (0,-7.75) {$n=n+1$};
  % 回边：向下、向左、向上，再水平进入判断框左顶点
  \draw[->] (-1.10,-7.75) -- (-3.45,-7.75) -- (-3.45,-4.55) -- (-2.55,-4.55);
  % 否：向右输出
  \draw[->] (2.55,-4.55) -- (3.95,-4.55);
  \node at (3.20,-4.20) {否};
  \draw (3.95,-4.10) -- (6.35,-4.10) -- (6.15,-5.00) -- (3.75,-5.00) -- cycle;
  \node at (5.05,-4.55) {输出 $n$};
  \draw[->] (5.05,-5.00) -- (5.05,-5.45);
  % 结束
  \draw[rounded corners=6pt] (4.10,-5.45) rectangle (6.00,-6.29);
  \node at (5.05,-5.87) {结束};
\end{tikzpicture}
